\section{信息增益：从提问到实验设计}
\label{sec:07-Information-Theory/06-Information-and-Learning/01}

\subsection{一个好问题值多少比特}

你在玩猜动物游戏，目标是确定一个隐藏类别 \(Y\)。可以问``它生活在水中吗''，也可以问``它名字的第二个字母是 a 吗''。两个问题都只有是/非两种答案，但直觉告诉你前者更有用——前者把动物界劈成了大致均匀的两个集合，后者几乎劈不动任何东西。但直觉不可靠——你需要一个数字来比较这两个问题的信息产出。

这个直觉的能量单位是熵。问之前，目标的不确定性为 \(H(Y)\)——如果 16 种动物等可能，\(H(Y)=4\) bit。一个理想的 yes/no 问题劈开可能性空间，减少恰好 1 bit 的不确定。但现实观测 \(X\) 不会每次都劈得利索：有时答案几乎敲定了 \(Y\)（不确定骤降）；有时答案模棱两可，和没问差不多。

决定在观测发生之前。你衡量的是\textbf{期望}——所有可能答案下剩余不确定度的加权平均。观察到 \(X=x\) 后，关于 \(Y\) 的剩余不确定性是 \(H(Y\mid X=x)\)；按 \(X\) 的分布取平均：

\[
H(Y\mid X)=\sum_x p(x)\,H(Y\mid X=x).
\]

二者之差即为信息增益，恰好等于互信息：

\[
\operatorname{IG}(Y;X)
=H(Y)-H(Y\mid X)
=I(Y;X).
\]

这里有一个容易被跳过的细节：期望算的是\textbf{看到答案之前}对这次提问价值的预估。你花钱安排实验，买的是期望信息增益——不是赌某次结果特别走运。就像你不会因为一次彩票中奖就论证买彩票是好策略。

信息增益的范围：下限为零（\(X\) 与 \(Y\) 独立，观测不做任何事），上限为 \(H(Y)\)（\(X\) 完全揭示 \(Y\)）。绝大多数真实观测落在这两个极端之间——提供一些信息，但不完整。

\subsection{一次决策树划分怎样计算：两个分裂的优劣一目了然}

某叶节点中有 8 个样本，正负类各 4 个，父节点熵为 1 bit。现在评估两个候选划分。

候选划分 \(A\)：把 4 个正例和 4 个负例完全分开——左子节点全正，右子节点全负。两个子节点纯净，条件熵为零：

\[
\operatorname{IG}_A=1-0=1\text{ bit}.
\]

候选划分 \(B\)：左子节点含 \((3+,1-)\)，右子节点含 \((1+,3-)\)。每个子节点的熵都是二元熵 \(h_2(3/4)\approx0.811\) bit，各占一半权重：

\[
\operatorname{IG}_B
=1-\frac12h_2(3/4)-\frac12h_2(3/4)
\approx0.189\text{ bit}.
\]

\(A\) 完胜。两个都是 yes/no 分裂，但一个把不确定度清零，另一个只削掉约五分之一。信息增益自动把``劈得干不干净''转成可排序的数字——决策树在每个节点只要算所有候选特征的 IG，选最大的。

但坑也埋在排序里。划分 \(C\)——按样本 ID 分裂，8 个样本各进一个分支，全部纯净——经验 IG 同样是 1 bit，和完美的二路分裂 \(A\) 同分。多路分裂天然占便宜：``切得碎''被奖励，即使碎片内部的纯净来自样本量太少而非真实规律。C4.5 用 gain ratio（IG 除以分裂本身的信息量）惩罚多路：8 路分裂需要的``分裂信息''是 \(h_2(8)=3\) bit，gain ratio 自动给 ID 分裂打出约 0.33 的分数，远低于完美二路分裂的 1.0。

另一个方向的选择偏差：如果你比较 500 个候选特征，从中挑 IG 最大者分裂——最大值天然高于单个真实 IG。就像 500 个人分别抛硬币，最优者可能抛出 7 次正面中的 7 次，但换成新硬币后会立刻回归均值。这种选择后的 IG 下降在决策树里对应``训练集上分裂好、测试集上不分裂更好''的典型过拟合信号。

\subsection{后验远离先验的那一刻，才叫学到了东西}

互信息还有一个对学习更直接的写法——把实验前和实验后的信念变化写成 KL 散度：

\[
I(Y;X)
=\E_X\left[
D\bigl(P_{Y\mid X}\|P_Y\bigr)
\right].
\]

对某个具体结果 \(x\)，KL 散度 \(D(P_{Y\mid x}\|P_Y)\) 衡量这次观察把你关于 \(Y\) 的信念推离了先验多远。如果看到 \(x\) 后，后验几乎等于先验（KL 接近零），观察近乎浪费。如果后验集中在单一类别而先验是 4 类等可能——KL = 2 bit——你一次学到了相当于原有全部不确定的信息量。

再按所有可能观察结果的出现概率加权平均，就是实验前对这次实验的合理估值。这个形式把三个看似独立的任务串在一起：

\begin{itemize}
\tightlist
\item \textbf{决策树}：\(X\) 是候选分裂特征，选 IG 最大。
\item \textbf{Bayesian 实验设计}：\(\Theta\) 是未知参数，\(D\) 是计划收集的数据。选实验条件 \(a\) 最大化
\[
I(\Theta;D\mid a)
=\E_{D\mid a}
D\bigl(P_{\Theta\mid D,a}\|P_\Theta\bigr).
\]
期望算在数据出现\textbf{之前}——实验前定价。
\item \textbf{主动学习}：你有标注预算，在未标注池中选``最值得标''的样本。标注后最能推离当前参数后验的样本——IG 最高——理应先标。
\end{itemize}

关键区分：事后 KL 是``这次实际学了多少''；事前期望是所有可能数据的平均——后者才是选实验的依据。一次好结果可能是运气，期望不靠单次胜利论证。

记号冲突：\(D(\cdot\|\cdot)\) 是 KL 散度，单独的 \(D\) 与 \(d\) 是观测数据。下一节最小描述长度也沿用 \(D\) 表示数据。

\subsubsection{两假设实验：完美信号与带噪信号}

两个假设 \(\Theta=0,1\) 先验等可能。实验一完美揭示真值——后验立刻变成点质量，KL = 1 bit，期望 IG = 1 bit。这是两个假设下的理论上限。

实验二以概率 0.75 报对、0.25 报错——等价于公平输入通过 BSC(0.25)。期望 IG = \(1-h_2(0.25)\approx0.189\) bit。一次带噪实验只值约五分之一个完美实验。想用 BSC(0.25) 获得接近 1 bit 的信息，大约要重复五六次，靠多数投票或贝叶斯更新逐步累积。这也恰好是 BSC(0.25) 的容量——先验刚好是最优分布。与决策树划分 \(B\) 给出相同的 0.189 bit 并非巧合：\(h_2(0.25)=h_2(0.75)\)，两处算的都是``四分之三可靠、四分之一随机''的同一个二元熵。

\subsection{新特征的价值取决于你已经知道什么}

若已有上下文 \(Z\)，新增特征 \(X\) 的价值写成条件形式：

\[
I(Y;X\mid Z).
\]

链式法则

\[
I(Y;X,Z)=I(Y;Z)+I(Y;X\mid Z)
\]

讲清了一件特征选择中反复被忽略的关键事实。``身高/cm''和``身高/inch''单独看都有可观的 IG——但在已知其中一个后，另一个的条件 IG 趋近于零。单独 IG 高不构成入选的充分理由——需要看\textbf{新增}信息。

更反直觉的是反面：单独 IG 为零的两个特征，联合可能完美预测标签。经典 XOR：\(X_1,X_2\) 独立公平比特，\(Y=X_1\oplus X_2\)。

单变量计算：\(I(Y;X_1)=I(Y;X_2)=0\)。知道 \(X_1=0\) 后，\(Y\) 仍以各半概率取 0 或 1——条件分布等于边缘分布，KL = 0，严格无信息。

联合计算：\(I(Y;X_1,X_2)=1\) bit。\((X_1,X_2)=(0,0)\) 或 \((1,1)\) 时 \(Y=0\)；\((0,1)\) 或 \((1,0)\) 时 \(Y=1\)。对每组 \((x_1,x_2)\)，\(Y\) 条件熵为零——联合预测确定。

单变量 IG 排序判定 \(X_1\) 和 \(X_2\) 都无价值，双双丢弃，你永远发现不了 XOR。特征选择不是排一列边际 IG 然后取前 k——某些特征的价值是\textbf{协同}的，只在与其他特征共存时显现。

推而广之，三个变量 \(X_1,X_2,X_3\) 时可能出现三层结构：\(X_1\) 单独有信息；\(X_2\) 单独无信息但跟 \(X_1\) 合用时补强；\(X_3\) 单独无信息、跟任一单变量合用也无信息、但三者联合才解锁。这种层级协同在基因调控网络和高阶交互中常见，边际 IG 排序完全看不见。

\subsection{从医学实验设计看信息增益的实际使用}

假设你怀疑一种新的降压药有疗效，设计了两种可能的实验来检验它。

\textbf{实验方案 A}：随机分两组（用药 vs 安慰剂），每组 50 人，测量收缩压变化。你预期的效应量约为 5 mmHg，个体间标准差约 10 mmHg。功效分析告诉你，这个设计有约 80\% 概率检测到显著差异（若效应真实存在）。

\textbf{实验方案 B}：同样的样本总量，但把 100 人按基线血压高低分成四层，每层内随机分配——控制了基线血压这个已知强预测因子。每层内方差减小，同样样本量下对效应量的估计更精确。

你用信息增益怎么比较 A 和 B？把未知效应量看作随机变量 \(\Theta\)，观测数据 \(D_A\) 或 \(D_B\) 是你收集的血压测量。B 方案减少了观测噪声（分层排除了一个已知变异源），所以 \(I(\Theta;D_B)>I(\Theta;D_A)\)——方案 B 在同样的样本成本下买了更多关于效应量的信息。实验设计的``功效''和``信息增益''在此等价：B 的功效更高，IG 也更高。

如果预算只够做一次实验，你还需要考虑分层测量本身的成本——也许分层需要额外的筛选检测。此时优化的不是 IG 本身，而是 IG 除以成本（每 bit 信息的价格）。最佳实验设计是在预算约束下最大化期望信息增益——不一定是 IG 绝对值最高的那个。

\subsection{经验信息增益为什么会对你说谎}

总体分布下独立变量的互信息为零——理论如此。有限样本中，用经验频率代入公式算出的 IG 几乎总是正数。本该是零的地方，经验 IG 也报出一个响亮的小正数（0.02 bit）。

偏差在特征数多时被选择效应急剧放大。从 5000 个候选逐个算 IG、挑最大——你看到的是``真实信号 + 抽样噪声''的最大值。最大值天然高于单次期望——5000 人同时各抛 10 枚硬币，总有人抛出 9 正面。你指着那人说硬币有偏，结论立足选择偏差而非物理。

最极致：高基数标识符。按 ID 分裂让每个叶节点只含一个样本，经验 IG = \(\log_2 n\)（1000 样本约 10 bit）——纯靠背训练集拿到满分。缓解方案各有代价：限制叶节点大小直接砍掉小分支；gain ratio 压低多路分裂权重；置换检验打乱特征值看 IG 下降是否显著；留出验证集评估真实泛化 IG。但每种都在暗中改变了选择准则——惩罚后的 IG 不再等于原来的期望信息量。

更深陷阱：IG 只度量\textbf{关联}，不度量\textbf{因果}。医院编号与治愈率高度相关——可能是三甲收治更多重症，不是挂哪家医院改病情。高 IG 特征未必是可干预杠杆。实验设计中尤为致命：你想最大化的是因果效应参数的 IG，不是被混杂污染的代理变量的 IG。按后者设计的实验，数据再多也推不出因果——信息增益的度量对象选错了。
